\documentclass[14pt, a4paper]{extarticle}
\usepackage[T2A]{fontenc}
\usepackage[utf8]{inputenc}
\usepackage[english,russian]{babel}
\usepackage[left = 3cm, right = 1cm, top = 2cm, bottom = 2cm]{geometry}
\setlength\parindent{12.5mm} 
\linespread{1.3}
\usepackage{enumitem}
\usepackage{amsmath}
\usepackage{amssymb}
\usepackage{wrapfig}
\usepackage{indentfirst}
\usepackage{hyperref}
\hypersetup{
	colorlinks=true,
	citecolor=blue,
	urlcolor=blue, 
	linkcolor=blue,
	pdfencoding=auto
}
\usepackage{xcolor}
\usepackage{listings}
\usepackage{graphicx}
\usepackage{float}
\usepackage{longtable}
\usepackage[htt]{hyphenat}
\usepackage{tabularx}

\definecolor{codegreen}{rgb}{0,0.6,0}
\definecolor{codegray}{rgb}{0.5,0.5,0.5}
\definecolor{codepurple}{rgb}{0.58,0,0.82}
\definecolor{backcolour}{rgb}{0.97,0.97,0.97}

\lstset{ 
	language=C++,
	backgroundcolor=\color{backcolour},
	commentstyle=\color{codegreen}\itshape,
	keywordstyle=\color{blue}\bfseries,
	numberstyle=\tiny\color{codegray},
	stringstyle=\color{codepurple},
	basicstyle=\ttfamily\footnotesize,
	breakatwhitespace=false,         
	breaklines=true,
	captionpos=t,
	keepspaces=true,                 
	numbers=left,
	numbersep=10pt,
	showspaces=false,                
	showstringspaces=false,
	showtabs=false,                  
	tabsize=4,
	frame=single,
	rulecolor=\color{black!30},
	extendedchars=\true
}

\newcommand{\tab}{\hspace{12.5mm}}
\newcolumntype{Y}{>{\centering\arraybackslash}X}

\usepackage{colortbl}

\definecolor{opcol}{RGB}{200, 240, 200}
\definecolor{modcol}{RGB}{200, 225, 250}
\definecolor{sibcol}{RGB}{255, 245, 190}
\definecolor{dispcol}{RGB}{255, 215, 215}
\definecolor{immcol}{RGB}{235, 215, 255}
\definecolor{rexcol}{RGB}{240, 230, 210}

\newcommand{\bytebox}[3]{%
\begin{tabular}{|c|c|c|c|c|c|c|c|}
\hline
\rowcolor{#1}
{\tiny 7} & {\tiny 6} & {\tiny 5} & {\tiny 4} & {\tiny 3} & {\tiny 2} & {\tiny 1} & {\tiny 0} \\ \hline
#2 \\ \hline
\multicolumn{8}{|c|}{\scriptsize #3} \\ \hline
\end{tabular}%
}

\newcommand{\modrmbox}[5]{%
\begin{tabular}{|c|c|c|c|c|c|c|c|}
\hline
\rowcolor{modcol}
{\tiny 7} & {\tiny 6} & {\tiny 5} & {\tiny 4} & {\tiny 3} & {\tiny 2} & {\tiny 1} & {\tiny 0} \\ \hline
#1 \\ \hline
\multicolumn{2}{|c|}{\tiny #2} & \multicolumn{3}{c|}{\tiny #3} & \multicolumn{3}{c|}{\tiny #4} \\ \hline
\multicolumn{8}{|c|}{\scriptsize #5} \\ \hline
\end{tabular}%
}

\newcommand{\sibbox}[5]{%
\begin{tabular}{|c|c|c|c|c|c|c|c|}
\hline
\rowcolor{sibcol}
{\tiny 7} & {\tiny 6} & {\tiny 5} & {\tiny 4} & {\tiny 3} & {\tiny 2} & {\tiny 1} & {\tiny 0} \\ \hline
#1 \\ \hline
\multicolumn{2}{|c|}{\tiny #2} & \multicolumn{3}{c|}{\tiny #3} & \multicolumn{3}{c|}{\tiny #4} \\ \hline
\multicolumn{8}{|c|}{\scriptsize #5} \\ \hline
\end{tabular}%
}

\newcommand{\rexbox}[6]{%
\begin{tabular}{|c|c|c|c|c|c|c|c|}
\hline
\rowcolor{rexcol}
{\tiny 7} & {\tiny 6} & {\tiny 5} & {\tiny 4} & {\tiny 3} & {\tiny 2} & {\tiny 1} & {\tiny 0} \\ \hline
#1 \\ \hline
\multicolumn{4}{|c|}{\tiny 0100 (REX)} & {\tiny #2} & {\tiny #3} & {\tiny #4} & {\tiny #5} \\ \hline
\multicolumn{8}{|c|}{\scriptsize #6} \\ \hline
\end{tabular}%
}

\begin{document}
	\thispagestyle{empty}
	\begin{center}
		\hspace{0.2cm}МИНИСТЕРСТВО НАУКИ И ВЫСШЕГО ОБРАЗОВАНИЯ РОССИЙСКОЙ ФЕДЕРАЦИИ\\
		Федеральное государственное автономное образовательное\\
		учреждение высшего образования\\
		{\bf <<Санкт-Петербургский политехнический\\
			университет Петра Великого>>}\\
		Институт компьютерных наук и кибербезопасности\\
		Высшая школа технологий искусственного интеллекта\\
		02.03.01\_03 Технологии машинного обучения и искусственного интеллекта\\
		\vspace{3cm}
		\Large\bf{
			Отчёт о выполнении лабораторной работы №6\\
			Основы архитектуры ЦВМ\\[1cm]
		}
	\end{center}
	\vspace{2cm}
	\begin{flushleft}
		{\bf Выполнил:}\\
		Cтудент группы 5130201/50302 \hfill
		\line(1,0){4cm} \quad Калинский В.И.\\
		{\bf Принял:}\\
		К.т.н., доцент, доцент ВШТИИ и КОЦ ИИ \hfill
		\line(1,0){3 cm} \quad Вербова Н.М.\\[0.5cm]
		{ \hfill<<\rule{1cm}{0.5pt}>>\rule{2cm}{0.5pt}\hspace{3pt}2026 г.}
	\end{flushleft}
	\vspace{1cm}
	\begin{center}
		Санкт-Петербург, 2026
	\end{center}
	\newpage

	\noindent \textbf{Цель работы:}
	
	Познакомиться с системой процессорных команд и способами использования этих команд. Освоить методику изучения регистровой модели, системы команд и способов адресации с использованием экранного отладчика среды Visual Studio.
	\\\\
	\noindent \textbf{Постановка задачи (вариант 2-8):}
	
	Для группы 50302 в соответствии с методическими указаниями задан вариант 2-8: массив \texttt{int iMas[10]}, итератор \texttt{unsigned int i}. В цикле \texttt{for (i = 0; i < 9; i++)} реализуется вычисление элементов массива: если $i \ge 4$, то $iMas[i] = 13 \cdot i + 7$, иначе $iMas[i] = (i \cdot 8) \oplus 12$. В подготовительной части исследуется прямое и косвенное присваивание констант различной длины: 12, 12345, 12345678, \texttt{0x123456789ABCDE} и \texttt{0xAB12} через указатель \texttt{pI = \&iMas[2]}.
	\\\\
	\noindent \textbf{1. Исходная программа и сопоставление с дизассемблером}
	
	В соответствии с шагом 1 методики в программе-шаблоне циклический фрагмент заменён кодом варианта 2-8 (листинг 1).
	
	\begin{lstlisting}[caption={Исходная программа (вариант 2-8, глобальные переменные)}]
int iMas[10];
int *pI;
unsigned int i;

int main() {
	iMas[3] = 12;
	iMas[4] = 12345;
	iMas[5] = 12345678;
	iMas[6] = 0x123456789ABCDE;

	pI = &iMas[2];
	*pI = 0xAB12;

	for (i = 0; i < 9; i++) {
		if (i >= 4)
			iMas[i] = 13 * i + 7;
		else
			iMas[i] = (i * 8) ^ 0xC;
	}
	return 0;
}\end{lstlisting}

	Программа оттранслирована в среде Visual Studio. В режиме отладки (окно Disassembly) получен машинный код (листинг 2) с сопоставлением инструкциям языка Си.

	\begin{lstlisting}[language={[x86masm]Assembler}, caption={Дизассемблированный код с сопоставлением инструкциям Си}]
     1: int iMas[10];
2: int *pI;
3: unsigned int i;
4: 
5: int main() {
	00007FF630CD1720 40 55                push        rbp                                  ; Сохранение базового регистра кадра
	00007FF630CD1722 57                   push        rdi                                  ; Сохранение регистра rdi
	00007FF630CD1723 48 81 EC E8 00 00 00 sub         rsp,0E8h                             ; Выделение памяти в стеке под локальные нужды
	00007FF630CD172A 48 8D 6C 24 20       lea         rbp,[rsp+20h]                        ; Установка базового указателя кадра
	00007FF630CD172F 48 8D 0D CA F8 00 00 lea         rcx,[__B952A550_main@cpp (07FF630CE1000h)] ; Служебный вызов отладчика
	00007FF630CD1736 E8 3A FC FF FF       call        __CheckForDebuggerJustMyCode (07FF630CD1375h) ; Проверка режима Just My Code
	00007FF630CD173B 90                   nop                                              ; Выравнивание
	6: 	iMas[3] = 12;
	00007FF630CD173C B8 04 00 00 00       mov         eax,4                                ; Размер элемента (4 байта для int)
	00007FF630CD1741 48 6B C0 03          imul        rax,rax,3                            ; Смещение элемента: 3 * 4 = 12 (0x0C) байт
	00007FF630CD1745 48 8D 0D F4 AA 00 00 lea         rcx,[iMas (07FF630CDC240h)]          ; Загрузка адреса массива в rcx
	00007FF630CD174C C7 04 01 0C 00 00 00 mov         dword ptr [rcx+rax],0Ch              ; Базово-индексная запись константы 12
	7: 	iMas[4] = 12345;
	00007FF630CD1753 B8 04 00 00 00       mov         eax,4                                ; Размер элемента
	00007FF630CD1758 48 6B C0 04          imul        rax,rax,4                            ; Смещение элемента: 4 * 4 = 16 (0x10) байт
	00007FF630CD175C 48 8D 0D DD AA 00 00 lea         rcx,[iMas (07FF630CDC240h)]          ; Базовый адрес массива iMas
	00007FF630CD1763 C7 04 01 39 30 00 00 mov         dword ptr [rcx+rax],3039h            ; Запись 2-байтовой константы 12345
	8: 	iMas[5] = 12345678;
	00007FF630CD176A B8 04 00 00 00       mov         eax,4                                ; Размер элемента
	00007FF630CD176F 48 6B C0 05          imul        rax,rax,5                            ; Смещение элемента: 5 * 4 = 20 (0x14) байт
	00007FF630CD1773 48 8D 0D C6 AA 00 00 lea         rcx,[iMas (07FF630CDC240h)]          ; Базовый адрес массива iMas
	00007FF630CD177A C7 04 01 4E 61 BC 00 mov         dword ptr [rcx+rax],0BC614Eh         ; Запись 4-байтовой константы 12345678
	9: 	iMas[6] = 0x123456789ABCDE;
	00007FF630CD1781 B8 04 00 00 00       mov         eax,4                                ; Размер элемента
	00007FF630CD1786 48 6B C0 06          imul        rax,rax,6                            ; Смещение элемента: 6 * 4 = 24 (0x18) байта
	00007FF630CD178A 48 8D 0D AF AA 00 00 lea         rcx,[iMas (07FF630CDC240h)]          ; Базовый адрес массива iMas
	00007FF630CD1791 C7 04 01 DE BC 9A 78 mov         dword ptr [rcx+rax],789ABCDEh        ; Запись младших 32 бит (усечение константы)
	10: 
	11: 	pI = &iMas[2];
	00007FF630CD1798 B8 04 00 00 00       mov         eax,4                                ; Размер элемента
	00007FF630CD179D 48 6B C0 02          imul        rax,rax,2                            ; Смещение для индекса 2: 2 * 4 = 8 байт
	00007FF630CD17A1 48 8D 0D 98 AA 00 00 lea         rcx,[iMas (07FF630CDC240h)]          ; Базовый адрес массива iMas
	00007FF630CD17A8 48 03 C8             add         rcx,rax                              ; rcx = адрес элемента iMas[2]
	00007FF630CD17AB 48 8B C1             mov         rax,rcx                              ; Перенос адреса в rax
	00007FF630CD17AE 48 89 05 B3 AA 00 00 mov         qword ptr [pI (07FF630CDC268h)],rax  ; Сохранение адреса в указатель pI (8 байт)
	12: 	*pI = 0xAB12;
	00007FF630CD17B5 48 8B 05 AC AA 00 00 mov         rax,qword ptr [pI (07FF630CDC268h)]  ; Считывание адреса из указателя pI в rax
	00007FF630CD17BC C7 00 12 AB 00 00    mov         dword ptr [rax],0AB12h               ; Косвенная запись константы 0xAB12 по адресу в rax
	13: 
	14: 	for (i = 0; i < 9; i++) {
		00007FF630CD17C2 C7 05 A4 AA 00 00 00 00 00 00 mov         dword ptr [i (07FF630CDC270h)],0 ; Инициализация счетчика цикла: i = 0
		00007FF630CD17CC EB 0E                jmp         main+0BCh (07FF630CD17DCh)           ; Безусловный переход к проверке условия цикла
		00007FF630CD17CE 8B 05 9C AA 00 00    mov         eax,dword ptr [i (07FF630CDC270h)]   ; Загрузка значения i в eax
		00007FF630CD17D4 FF C0                inc         eax                                  ; Инкремент счетчика (i++)
		00007FF630CD17D6 89 05 94 AA 00 00    mov         dword ptr [i (07FF630CDC270h)],eax   ; Сохранение обновленного i
		00007FF630CD17DC 83 3D 8D AA 00 00 09 cmp         dword ptr [i (07FF630CDC270h)],9     ; Сравнение i с 9
		00007FF630CD17E3 73 43                jae         main+108h (07FF630CD1828h)           ; Если i >= 9, выход из цикла (беззнаковый переход)
		15: 		if (i >= 4)
		00007FF630CD17E5 83 3D 84 AA 00 00 04 cmp         dword ptr [i (07FF630CDC270h)],4     ; Сравнение i с 4
		00007FF630CD17EC 72 1C                jb          main+0EAh (07FF630CD180Ah)           ; Если i < 4, переход в ветку else (jb, CF=1)
		16: 			iMas[i] = 13 * i + 7;
		00007FF630CD17EE 6B 05 7B AA 00 00 0D imul        eax,dword ptr [i (07FF630CDC270h)],0Dh ; eax = i * 13 (0x0D)
		00007FF630CD17F5 83 C0 07             add         eax,7                                ; eax = 13 * i + 7
		00007FF630CD17F8 8B 0D 72 AA 00 00    mov         ecx,dword ptr [i (07FF630CDC270h)]   ; Загрузка индекса i в ecx
		00007FF630CD17FE 48 8D 15 3B AA 00 00 lea         rdx,[iMas (07FF630CDC240h)]          ; Базовый адрес массива в rdx
		00007FF630CD1805 89 04 8A             mov         dword ptr [rdx+rcx*4],eax            ; Запись в массив: База + Индекс * 4 (код 89 04 8A)
		00007FF630CD1808 EB 1C                jmp         main+106h (07FF630CD1826h)           ; Обход ветки else
		17: 		else
		18: 			iMas[i] = (i * 8) ^ 0xC;
		00007FF630CD180A 8B 05 60 AA 00 00    mov         eax,dword ptr [i (07FF630CDC270h)]   ; Загрузка i в eax
		00007FF630CD1810 C1 E0 03             shl         eax,3                                ; Умножение на 8 сдвигом влево на 3 бита
		00007FF630CD1813 83 F0 0C             xor         eax,0Ch                              ; Побитовое исключающее ИЛИ с константой 12
		00007FF630CD1816 8B 0D 54 AA 00 00    mov         ecx,dword ptr [i (07FF630CDC270h)]   ; Загрузка индекса i в ecx
		00007FF630CD181C 48 8D 15 1D AA 00 00 lea         rdx,[iMas (07FF630CDC240h)]          ; Базовый адрес массива в rdx
		00007FF630CD1823 89 04 8A             mov         dword ptr [rdx+rcx*4],eax            ; Запись в массив: База + Индекс * 4 (код 89 04 8A)
		19: 	}
	00007FF630CD1826 EB A6                jmp         main+0AEh (07FF630CD17CEh)           ; Переход к следующей итерации (на шаг цикла)
	20: 
	21: 	return 0;
	00007FF630CD1828 33 C0                xor         eax,eax                              ; Код возврата функции 0 (return 0)
	22: }
00007FF630CD182A 48 8D A5 C8 00 00 00 lea         rsp,[rbp+0C8h]                       ; Восстановление указателя стека
00007FF630CD1831 5F                   pop         rdi                                  ; Восстановление регистра rdi
00007FF630CD1832 5D                   pop         rbp                                  ; Восстановление базового указателя кадра rbp
00007FF630CD1833 C3                   ret                                              ; Возврат из функции
\end{lstlisting}
     	
	\noindent \textbf{2. Список переменных и карта памяти}
	
	Переменные \texttt{iMas}, \texttt{pI}, \texttt{i} объявлены вне функций (глобально) и размещены компилятором в статической памяти данных. Их адреса постоянны и зафиксированы отладчиком (таблица \ref{tab:mem_map}).
	
	\begin{table}[H]
		\centering
		\caption{Список переменных и карта памяти программы} \label{tab:mem_map}
		\begin{tabularx}{\textwidth}{|l|c|c|c|X|}
			\hline
			\textbf{Имя} & \textbf{Класс} & \textbf{Тип} & \textbf{Размер} & \textbf{Адрес в памяти} \\ \hline
			\texttt{iMas[0..9]} & глобальная & \texttt{int[10]} & 40 байт & \texttt{0x07FF630CDC240} -- \texttt{0x07FF630CDC267} \\ \hline
			\texttt{pI} & глобальная & \texttt{int*} & 8 байт & \texttt{0x07FF630CDC268} -- \texttt{0x07FF630CDC26F} \\ \hline
			\texttt{i} & глобальная & \texttt{unsigned int} & 4 байта & \texttt{0x07FF630CDC270} -- \texttt{0x07FF630CDC273} \\ \hline
		\end{tabularx}
	\end{table}

	Переменные расположены последовательно: массив \texttt{iMas} занимает 40 байт ($28_{16}$), следом 8-байтовый указатель \texttt{pI}, затем 4-байтовый счетчик \texttt{i}.
	\\\\
	\noindent \textbf{3. Анализ команд обращения к переменным}
	
	В таблице \ref{tab:instructions_analysis} систематизированы все команды, обращающиеся к переменным программы, с указанием числа операндов и способов их адресации, а в таблице \ref{tab:machine_commands_10} -- 10 команд байтовых в ассемблере с расшифровкой.
	
	\begin{table}[H]
		\centering
		\caption{Команды обращения к переменным программы} \label{tab:instructions_analysis}
		\begin{tabularx}{\textwidth}{|p{5.2cm}|c|c|X|X|}
			\hline
			\textbf{Команда ассемблера} & \textbf{Перем.} & \textbf{Оп.} & \textbf{Операнд 1 (приёмник)} & \textbf{Операнд 2 (источник)} \\ \hline
			\texttt{mov [rcx+rax], 0Ch} & \texttt{iMas[3]} & 2 & Базово-индексная & Непосредств. \\ \hline
			\texttt{mov [rcx+rax], 3039h} & \texttt{iMas[4]} & 2 & Базово-индексная & Непосредств. \\ \hline
			\texttt{mov [rcx+rax], 0BC614Eh} & \texttt{iMas[5]} & 2 & Базово-индексная & Непосредств. \\ \hline
			\texttt{mov [rcx+rax], 789ABCDEh} & \texttt{iMas[6]} & 2 & Базово-индексная & Непосредств. \\ \hline
			\texttt{mov qword ptr [pI], rax} & \texttt{pI} & 2 & Прямая & Регистровая \\ \hline
			\texttt{mov rax, qword ptr [pI]} & \texttt{pI} & 2 & Регистровая & Прямая \\ \hline
			\texttt{mov dword ptr [rax], 0AB12h} & \texttt{iMas[2]} & 2 & Косвенно-регистровая & Непосредств. \\ \hline
			\texttt{mov dword ptr [i], 0} & \texttt{i} & 2 & Прямая & Непосредств. \\ \hline
			\texttt{mov eax, dword ptr [i]} & \texttt{i} & 2 & Регистровая & Прямая \\ \hline
			\texttt{mov dword ptr [i], eax} & \texttt{i} & 2 & Прямая & Регистровая \\ \hline
			\texttt{cmp dword ptr [i], 9} & \texttt{i} & 2 & Прямая & Непосредств. \\ \hline
			\texttt{cmp dword ptr [i], 4} & \texttt{i} & 2 & Прямая & Непосредств. \\ \hline
			\texttt{imul eax, dword ptr [i], 0Dh} & \texttt{i} & 3 & Регистровая & Прямая и непоср. \\ \hline
			\texttt{mov ecx, dword ptr [i]} & \texttt{i} & 2 & Регистровая & Прямая \\ \hline
			\texttt{mov [rdx+rcx*4], eax} & \texttt{iMas[i]} & 2 & Масштабно-индексная & Регистровая \\ \hline
		\end{tabularx}
	\end{table}


	\noindent \textbf{4. Рассмотрим каждую ассемблерную команду подробно (процессор x86)}\\[3mm]
	Для 10 команд программы, систематизированных в таблице \ref{tab:machine_commands_10}, ниже представлен подробный побайтовый и побитовый разбор структуры машинного кода в соответствии с общим форматом команды архитектуры x86 (префиксы, код операции, байт ModR/M, байт SIB, поле смещения Displacement и непосредственный операнд Immediate).

\begin{table}[H]
	\centering
	\caption{Примеры машинных команд и их байтового представления} \label{tab:machine_commands_10}
	\begin{tabularx}{\textwidth}{|p{4.2cm}|p{4.2cm}|X|}
		\hline
		\textbf{Байт} & \textbf{Ассемблер} & \textbf{Пояснение} \\ \hline
		\texttt{40 55} & \texttt{push rbp} & Сохранение регистра rbp \\ \hline
		\texttt{48 81 EC E8 00 00 00} & \texttt{sub rsp, 0E8h} & Выделение стека \\ \hline
		\texttt{B8 04 00 00 00} & \texttt{mov eax, 4} & Загрузка 4 в eax \\ \hline
		\texttt{48 6B C0 03} & \texttt{imul rax, rax, 3} & Умножение rax на 3 \\ \hline
		\texttt{48 8D 6C 24 20 } & \texttt{lea rbp,[rsp+20h] } & Загрузка адреса в rbp \\ \hline
		\texttt{FF C0} & \texttt{inc eax} & Инкремент eax \\ \hline
		\texttt{83 3D 8D AA 00 00 09} & \texttt{cmp dword ptr [i], 9} & Сравнение i с 9 \\ \hline
		\texttt{73 43} & \texttt{jae main+108h (07FF630CD1828hP)} & Беззнаковый переход \\ \hline
		\texttt{33 C0} & \texttt{xor eax, eax} & Обнуление eax \\ \hline
		\texttt{C3} & \texttt{ret} & Возврат из функции \\ \hline
	\end{tabularx}
\end{table}

	% --- 1. PUSH RBP ---
	\noindent \textbf{\textcolor{red}{00007FF630CD1720}} \quad \textbf{\textcolor{codegreen}{40 55}} \quad \texttt{push rbp}\\[1.5mm]
	\noindent \resizebox{0.52\textwidth}{!}{%
	\rexbox{{\tiny 0}&{\tiny 1}&{\tiny 0}&{\tiny 0}&{\tiny 0}&{\tiny 0}&{\tiny 0}&{\tiny 0}}{0}{0}{0}{0}{REX (40h)} \quad
	\bytebox{opcol}{{\tiny 0}&{\tiny 1}&{\tiny 0}&{\tiny 1}&{\tiny 0}&{\tiny 1}&{\tiny 0}&{\tiny 1}}{Opcode (55h)}
	}\\[1.5mm]
	\texttt{5xh - push}\\
	\texttt{x5h - rbp}\\
	Неявная адресация \\[4mm]

	% --- 2. SUB RSP, 0E8h ---
	\noindent \textbf{\textcolor{red}{00007FF630CD1723}} \quad \textbf{\textcolor{codegreen}{48 81 EC E8 00 00 00}} \quad \texttt{sub rsp, 0E8h}\\[1.5mm]
	\noindent \resizebox{\textwidth}{!}{%
	\rexbox{{\tiny 0}&{\tiny 1}&{\tiny 0}&{\tiny 0}&{\tiny 1}&{\tiny 0}&{\tiny 0}&{\tiny 0}}{W=1}{0}{0}{0}{REX.W (48h)} \quad
	\bytebox{opcol}{{\tiny 1}&{\tiny 0}&{\tiny 0}&{\tiny 0}&{\tiny 0}&{\tiny 0}&{\tiny 0}&{\tiny 1}}{Opcode (81h)} \quad
	\modrmbox{{\tiny 1}&{\tiny 1}&{\tiny 1}&{\tiny 0}&{\tiny 1}&{\tiny 1}&{\tiny 0}&{\tiny 0}}{MOD=3}{REG=5}{R/M=4}{mod-reg-r/m (ECh)}
	}
	\\\\	
	\noindent \resizebox{\textwidth}{!}{
	\bytebox{immcol}{{\tiny 1}&{\tiny 1}&{\tiny 1}&{\tiny 0}&{\tiny 1}&{\tiny 0}&{\tiny 0}&{\tiny 0}}{Imm0 (E8h)} \quad
	\bytebox{immcol}{{\tiny 0}&{\tiny 0}&{\tiny 0}&{\tiny 0}&{\tiny 0}&{\tiny 0}&{\tiny 0}&{\tiny 0}}{Imm1 (00h)} \quad
	\bytebox{immcol}{{\tiny 0}&{\tiny 0}&{\tiny 0}&{\tiny 0}&{\tiny 0}&{\tiny 0}&{\tiny 0}&{\tiny 0}}{Imm2 (00h)} \quad
	\bytebox{immcol}{{\tiny 0}&{\tiny 0}&{\tiny 0}&{\tiny 0}&{\tiny 0}&{\tiny 0}&{\tiny 0}&{\tiny 0}}{Imm3 (00h)}
	}\\[1.5mm]
	\texttt{81 - SUB}\\
	\texttt{MOD = 3 - регистровая адресация (R/M содержит номер регистра)}\\
	\texttt{REG = 5 код операции SUB}\\
	\texttt{R/M = 4 - rsp}\\
	\texttt{Immediate = 000000E8h (выделение 232 байт в стековом кадре)}\\
	Непосредственная адресация\\[4mm]

	% --- 3. MOV EAX, 4 ---
	\noindent \textbf{\textcolor{red}{00007FF630CD173C}} \quad \textbf{\textcolor{codegreen}{B8 04 00 00 00}} \quad \texttt{mov eax, 4}\\[1.5mm]
	\noindent \resizebox{0.65\textwidth}{!}{%
	\bytebox{opcol}{{\tiny 1}&{\tiny 0}&{\tiny 1}&{\tiny 1}&{\tiny 1}&{\tiny 0}&{\tiny 0}&{\tiny 0}}{Opcode (B8h)} \quad
	\bytebox{immcol}{{\tiny 0}&{\tiny 0}&{\tiny 0}&{\tiny 0}&{\tiny 0}&{\tiny 1}&{\tiny 0}&{\tiny 0}}{Imm0 (04h)} \quad
}
\\
	\noindent \resizebox{0.85\textwidth}{!}{
	\bytebox{immcol}{{\tiny 0}&{\tiny 0}&{\tiny 0}&{\tiny 0}&{\tiny 0}&{\tiny 0}&{\tiny 0}&{\tiny 0}}{Imm1 (00h)} \quad
	\bytebox{immcol}{{\tiny 0}&{\tiny 0}&{\tiny 0}&{\tiny 0}&{\tiny 0}&{\tiny 0}&{\tiny 0}&{\tiny 0}}{Imm2 (00h)} \quad
	\bytebox{immcol}{{\tiny 0}&{\tiny 0}&{\tiny 0}&{\tiny 0}&{\tiny 0}&{\tiny 0}&{\tiny 0}&{\tiny 0}}{Imm3 (00h)}
	}\\[1.5mm]
	\texttt{B8 - код mov}\\
	\texttt{reg = 0 - eax}\\
	\texttt{Immediate = 00000004h }\\
	непосредственная адресация\\[4mm]

	% --- 4. IMUL RAX, RAX, 3 ---
	\noindent \textbf{\textcolor{red}{00007FF630CD1741}} \quad \textbf{\textcolor{codegreen}{48 6B C0 03}} \quad \texttt{imul rax, rax, 3}\\[1.5mm]
	\noindent \resizebox{0.75\textwidth}{!}{%
	\rexbox{{\tiny 0}&{\tiny 1}&{\tiny 0}&{\tiny 0}&{\tiny 1}&{\tiny 0}&{\tiny 0}&{\tiny 0}}{W=1}{0}{0}{0}{REX.W (48h)} \quad
	\bytebox{opcol}{{\tiny 0}&{\tiny 1}&{\tiny 1}&{\tiny 0}&{\tiny 1}&{\tiny 0}&{\tiny 1}&{\tiny 1}}{Opcode (6Bh)} \quad
}

	\noindent \resizebox{0.75\textwidth}{!}{
	\modrmbox{{\tiny 1}&{\tiny 1}&{\tiny 0}&{\tiny 0}&{\tiny 0}&{\tiny 0}&{\tiny 0}&{\tiny 0}}{MOD=3}{REG=0}{R/M=0}{mod-reg-r/m (C0h)} \quad
	\bytebox{immcol}{{\tiny 0}&{\tiny 0}&{\tiny 0}&{\tiny 0}&{\tiny 0}&{\tiny 0}&{\tiny 1}&{\tiny 1}}{Imm (03h)}
	}\\[1.5mm]
	\texttt{6B - код imul}\\	
	\texttt{MOD = 3 - регистровая адресация}\\
	\texttt{REG = 0 - eax}\\
	\texttt{R/M = 0 - eax}\\
	\texttt{Immediate = 03h}\\
	регистровая и непосредственная адресация\\[4mm]

	% --- 5. LEA RBP, [RSP+20h] ---
	\noindent \textbf{\textcolor{red}{00007FF630CD172A}} \quad \textbf{\textcolor{codegreen}{48 8D 6C 24 20}} \quad \texttt{lea rbp, [rsp+20h]}\\[1.5mm]
	\noindent \resizebox{0.8\textwidth}{!}{%
	\rexbox{{\tiny 0}&{\tiny 1}&{\tiny 0}&{\tiny 0}&{\tiny 1}&{\tiny 0}&{\tiny 0}&{\tiny 0}}{W=1}{0}{0}{0}{REX.W (48h)} \quad
	\bytebox{opcol}{{\tiny 1}&{\tiny 0}&{\tiny 0}&{\tiny 0}&{\tiny 1}&{\tiny 1}&{\tiny 0}&{\tiny 1}}{Opcode (8Dh)} \quad
}

	\noindent \resizebox{\textwidth}{!}{
	\modrmbox{{\tiny 0}&{\tiny 1}&{\tiny 1}&{\tiny 0}&{\tiny 1}&{\tiny 1}&{\tiny 0}&{\tiny 0}}{MOD=1}{REG=5}{R/M=4}{mod-reg-r/m (6Ch)} \quad
	\sibbox{{\tiny 0}&{\tiny 0}&{\tiny 1}&{\tiny 0}&{\tiny 0}&{\tiny 1}&{\tiny 0}&{\tiny 0}}{Scale=0}{Index=4}{Base=4}{SIB (24h)} \quad
	\bytebox{dispcol}{{\tiny 0}&{\tiny 0}&{\tiny 1}&{\tiny 0}&{\tiny 0}&{\tiny 0}&{\tiny 0}&{\tiny 0}}{Disp (20h)}
	}\\[1.5mm]
	\texttt{8D - LEA}\\
	\texttt{MOD = 1 - адресация памяти со смещением disp8}\\
	\texttt{REG = 5 - rbp}\\
	\texttt{R/M = 4 - смотрим SIB}\\
	\texttt{Scale = 0 (x1), Index = 4, Base = 4}\\
	\texttt{Displacement = 20h (смещение +32 байта)}\\
	косвенная регистровая адресация со смещением\\[4mm]

	% --- 6. INC EAX ---
	\noindent \textbf{\textcolor{red}{00007FF630CD17D4}} \quad \textbf{\textcolor{codegreen}{FF C0}} \quad \texttt{inc eax}\\[1.5mm]
	\noindent \resizebox{0.52\textwidth}{!}{%
	\bytebox{opcol}{{\tiny 1}&{\tiny 1}&{\tiny 1}&{\tiny 1}&{\tiny 1}&{\tiny 1}&{\tiny 1}&{\tiny 1}}{Opcode (FFh)} \quad
	\modrmbox{{\tiny 1}&{\tiny 1}&{\tiny 0}&{\tiny 0}&{\tiny 0}&{\tiny 0}&{\tiny 0}&{\tiny 0}}{MOD=3}{REG=0}{R/M=0}{mod-reg-r/m (C0h)}
	}\\[1.5mm]
	\texttt{FF - INC}\\
	\texttt{MOD = 3 - регистровая адресация}\\
	\texttt{REG = 0}\\
	\texttt{R/M = 0 - eax}\\
	регистровая адресация\\[4mm]

	% --- 7. CMP DWORD PTR [I], 9 ---
	\noindent \textbf{\textcolor{red}{00007FF630CD17DC}} \quad \textbf{\textcolor{codegreen}{83 3D 8D AA 00 00 09}} \quad \texttt{cmp dword ptr [i], 9}\\[1.5mm]
	\noindent \resizebox{\textwidth}{!}{%
	\bytebox{opcol}{{\tiny 1}&{\tiny 0}&{\tiny 0}&{\tiny 0}&{\tiny 0}&{\tiny 0}&{\tiny 1}&{\tiny 1}}{Opcode (83h)} \quad
	\modrmbox{{\tiny 0}&{\tiny 0}&{\tiny 1}&{\tiny 1}&{\tiny 1}&{\tiny 1}&{\tiny 0}&{\tiny 1}}{MOD=0}{REG=7}{R/M=5}{mod-reg-r/m (3Dh)} \quad
	\bytebox{dispcol}{{\tiny 1}&{\tiny 0}&{\tiny 0}&{\tiny 0}&{\tiny 1}&{\tiny 1}&{\tiny 0}&{\tiny 1}}{Disp0 (8Dh)} \quad
}

\noindent \resizebox{\textwidth}{!}{
	\bytebox{dispcol}{{\tiny 1}&{\tiny 0}&{\tiny 1}&{\tiny 0}&{\tiny 1}&{\tiny 0}&{\tiny 1}&{\tiny 0}}{Disp1 (AAh)} \quad
	\bytebox{dispcol}{{\tiny 0}&{\tiny 0}&{\tiny 0}&{\tiny 0}&{\tiny 0}&{\tiny 0}&{\tiny 0}&{\tiny 0}}{Disp2 (00h)} \quad
	\bytebox{dispcol}{{\tiny 0}&{\tiny 0}&{\tiny 0}&{\tiny 0}&{\tiny 0}&{\tiny 0}&{\tiny 0}&{\tiny 0}}{Disp3 (00h)} \quad
	\bytebox{immcol}{{\tiny 0}&{\tiny 0}&{\tiny 0}&{\tiny 0}&{\tiny 1}&{\tiny 0}&{\tiny 0}&{\tiny 1}}{Immediate (09h)}
	}\\[1.5mm]
	\texttt{83 - CMP}\\
	\texttt{MOD = 0 - displacement содержит адрес переменной в памяти}\\
	\texttt{REG = 7 - CMP}\\
	\texttt{R/M = 5 ebx}\\
	\texttt{Displacement = 0000AA8Dh}\\
	\texttt{Immediate = 09h}\\
	прямая адресация памяти и непосредственная адресация\\[4mm]

	% --- 8. JAE ---
	\noindent \textbf{\textcolor{red}{00007FF630CD17E3}} \quad \textbf{\textcolor{codegreen}{73 43}} \quad \texttt{jae main+108h (07FF630CD1828h)}\\[1.5mm]
	\noindent \resizebox{0.52\textwidth}{!}{%
	\bytebox{opcol}{{\tiny 0}&{\tiny 1}&{\tiny 1}&{\tiny 1}&{\tiny 0}&{\tiny 0}&{\tiny 1}&{\tiny 1}}{Opcode (73h)} \quad
	\bytebox{dispcol}{{\tiny 0}&{\tiny 1}&{\tiny 0}&{\tiny 0}&{\tiny 0}&{\tiny 0}&{\tiny 1}&{\tiny 1}}{Displacement (43h)}
	}\\[1.5mm]
	\texttt{73 cb - JAE}\\
	\texttt{IP = 00007FF630CD17E5h (адрес следующей команды)}\\
	\texttt{next\_command\_address = IP + 0x43 = 07FF630CD1828h, если соблюдается условие ($i \ge 9$)}\\
	относительная адресация\\[4mm]

	% --- 9. XOR EAX, EAX ---
	\noindent \textbf{\textcolor{red}{00007FF630CD1828}} \quad \textbf{\textcolor{codegreen}{33 C0}} \quad \texttt{xor eax, eax}\\[1.5mm]
	\noindent \resizebox{0.52\textwidth}{!}{%
	\bytebox{opcol}{{\tiny 0}&{\tiny 0}&{\tiny 1}&{\tiny 1}&{\tiny 0}&{\tiny 0}&{\tiny 1}&{\tiny 1}}{Opcode (33h)} \quad
	\modrmbox{{\tiny 1}&{\tiny 1}&{\tiny 0}&{\tiny 0}&{\tiny 0}&{\tiny 0}&{\tiny 0}&{\tiny 0}}{MOD=3}{REG=0}{R/M=0}{mod-reg-r/m (C0h)}
	}\\[1.5mm]
	\texttt{33 - XOR}\\
	\texttt{MOD = 3 - регистровая адресация (R/M содержит номер регистра)}\\
	\texttt{REG = 0 - eax}\\
	\texttt{R/M = 0 - eax}\\
	регистровая адресация\\[4mm]

	% --- 10. RET ---
	\noindent \textbf{\textcolor{red}{00007FF630CD1833}} \quad \textbf{\textcolor{codegreen}{C3}} \quad \texttt{ret}\\[1.5mm]
	\noindent \resizebox{0.28\textwidth}{!}{%
	\bytebox{opcol}{{\tiny 1}&{\tiny 1}&{\tiny 0}&{\tiny 0}&{\tiny 0}&{\tiny 0}&{\tiny 1}&{\tiny 1}}{Opcode (C3h)}
	}\\[1.5mm]
	\texttt{C3 - RET}\\
	неявная адресация



	\noindent \textbf{5. Многокомпонентная адресация данных}
	
	Многокомпонентная масштабно-индексная адресация использована в команде записи в массив в теле цикла:
	\begin{equation*}
		\texttt{mov dword ptr [rdx+rcx*4], eax} \quad (\text{машинный код: } \texttt{89 04 8A})
	\end{equation*}
	Значения отдельных компонент исполнительного адреса:
	\begin{itemize}[noitemsep]
		\item \textbf{База (Base):} регистр \texttt{rdx} --- содержит базовый адрес начала массива \texttt{iMas} (\texttt{0x07FF630CDC240});
		\item \textbf{Индекс (Index):} регистр \texttt{rcx} --- содержит значение переменной \texttt{i} ($0 \le i < 9$);
		\item \textbf{Масштаб (Scale):} коэффициент 4 --- соответствует размеру типа \texttt{sizeof(int) = 4} байта;
		\item \textbf{Смещение (Displacement):} отсутствует (равно 0).
	\end{itemize}
	Исполнительный адрес вычисляется аппаратно за один такт: $EA = rdx + rcx \cdot 4$.
	\\\\
	\noindent \textbf{6. Способы адресации в командах перехода}
	
	Во всех командах ветвления компилятор применил \textbf{относительную адресацию} (\texttt{rel8}): адрес перехода не хранится целиком, а задаётся знаковым смещением относительно счётчика команд следующей инструкции ($\text{Target} = \text{IP} + \text{offset}$):
	\begin{itemize}[noitemsep]
		\item \texttt{EB 0E (jmp main+0BCh)}: безусловный короткий переход на проверку условия;
		\item \texttt{73 43 (jae main+108h)}: беззнаковый переход при выходе из цикла ($i \ge 9$);
		\item \texttt{72 1C (jb main+0EAh)}: беззнаковый переход в ветку \texttt{else} при $i < 4$ (по флагу $CF = 1$);
		\item \texttt{EB 1C (jmp main+106h)}: безусловный переход в обход ветки \texttt{else};
		\item \texttt{EB A6 (jmp main+0CEh)}: безусловный переход на шаг цикла со смещением \texttt{0xA6} ($-90$).
	\end{itemize}

	\noindent \textbf{7. Модификация программы: доступ через указатель и локальные переменные}
	
	В соответствии с шагом 10 методики доступ к массиву изменён на последовательное обращение через указатель \texttt{pMas} (листинг 3).
	
	\begin{lstlisting}[caption={Модифицированная программа с доступом через указатель}]
int iMas[10];
int *pI;
unsigned int i;

int main() {
	iMas[3] = 12;
	iMas[4] = 12345;
	iMas[5] = 12345678;
	iMas[6] = 0x123456789ABCDE;

	pI = &iMas[2];
	*pI = 0xAB12;

	int *pMas = iMas;
	for (i = 0; i < 9; i++, pMas++) {
		if (i >= 4)
			*pMas = 13 * i + 7;
		else
			*pMas = (i * 8) ^ 0xC;
	}
	return 0;
}
	\end{lstlisting}

	Дизассемблированный фрагмент модифицированного цикла:
	\begin{lstlisting}[language={[x86masm]Assembler}, caption={Дизассемблерный код цикла при работе через указатель}]
00007FF721F317C2 48 8D 05 77 AA 00 00 lea         rax,[iMas (07FF721F3C240h)]          ; Получение адреса начала массива iMas
00007FF721F317C9 48 89 45 08          mov         qword ptr [pMas],rax                 ; Сохранение адреса в локальный указатель pMas
15: 	for (i = 0; i < 9; i++, pMas++) {
	00007FF721F317CD C7 05 99 AA 00 00 00 00 00 00 mov         dword ptr [i (07FF721F3C270h)],0 ; Инициализация счетчика цикла: i = 0
	00007FF721F317D7 EB 1A                jmp         main+0D3h (07FF721F317F3h)           ; Безусловный переход к проверке условия цикла
	00007FF721F317D9 8B 05 91 AA 00 00    mov         eax,dword ptr [i (07FF721F3C270h)]   ; Загрузка значения i в eax
	00007FF721F317DF FF C0                inc         eax                                  ; Инкремент счетчика (i++)
	00007FF721F317E1 89 05 89 AA 00 00    mov         dword ptr [i (07FF721F3C270h)],eax   ; Сохранение обновленного i
	00007FF721F317E7 48 8B 45 08          mov         rax,qword ptr [pMas]                 ; Загрузка адреса текущего элемента из pMas
	00007FF721F317EB 48 83 C0 04          add         rax,4                                ; Инкремент указателя на 4 байта (pMas++)
	00007FF721F317EF 48 89 45 08          mov         qword ptr [pMas],rax                 ; Сохранение нового адреса в pMas
	00007FF721F317F3 83 3D 76 AA 00 00 09 cmp         dword ptr [i (07FF721F3C270h)],9     ; Сравнение i с 9
	00007FF721F317FA 73 2F                jae         main+10Bh (07FF721F3182Bh)           ; Если i >= 9, выход из цикла (беззнаковый переход)
	16: 		if (i >= 4)
	00007FF721F317FC 83 3D 6D AA 00 00 04 cmp         dword ptr [i (07FF721F3C270h)],4     ; Сравнение i с 4
	00007FF721F31803 72 12                jb          main+0F7h (07FF721F31817h)           ; Если i < 4, переход в ветку else (jb, CF=1)
	17: 			*pMas = 13 * i + 7;
	00007FF721F31805 6B 05 64 AA 00 00 0D imul        eax,dword ptr [i (07FF721F3C270h)],0Dh ; eax = i * 13 (0x0D)
	00007FF721F3180C 83 C0 07             add         eax,7                                ; eax = 13 * i + 7
	00007FF721F3180F 48 8B 4D 08          mov         rcx,qword ptr [pMas]                 ; Загрузка адреса элемента из указателя pMas в rcx
	00007FF721F31813 89 01                mov         dword ptr [rcx],eax                  ; Косвенно-регистровая запись [rcx] (код 89 01)
	00007FF721F31815 EB 12                jmp         main+109h (07FF721F31829h)           ; Обход ветки else
	18: 		else
	19: 			*pMas = (i * 8) ^ 0xC;
	00007FF721F31817 8B 05 53 AA 00 00    mov         eax,dword ptr [i (07FF721F3C270h)]   ; Загрузка i в eax
	00007FF721F3181D C1 E0 03             shl         eax,3                                ; Умножение на 8 сдвигом влево на 3 бита
	00007FF721F31820 83 F0 0C             xor         eax,0Ch                              ; Побитовое исключающее ИЛИ с константой 12
	00007FF721F31823 48 8B 4D 08          mov         rcx,qword ptr [pMas]                 ; Загрузка адреса элемента из указателя pMas в rcx
	00007FF721F31827 89 01                mov         dword ptr [rcx],eax                  ; Косвенно-регистровая запись [rcx] (код 89 01)\end{lstlisting}

	\noindent \textbf{Сравнение способов адресации:}
	\begin{enumerate}[noitemsep]
		\item \textbf{Индекс против указателя:} При доступе по индексу процессор на каждом шаге вычисляет формулу $rdx + rcx \cdot 4$ командой \texttt{mov [rdx+rcx*4], eax} (машинный код \texttt{89 04 8A}, 3 байта, содержит SIB-байт). При доступе через указатель запись выполняется короткой командой косвенно-регистровой адресации \texttt{mov [rcx], eax} (код \texttt{89 01}, 2 байта, без SIB-байта), а шаг к следующему элементу сводится к простому сложению \texttt{add rax, 4}.
		\item \textbf{Глобальное против локального размещения:} Глобальные переменные размещены в статической памяти и адресуются по фиксированным смещениям. Локальные переменные размещаются в стеке и адресуются со смещением относительно регистра базы кадра: \texttt{[rbp-offset]}.
	\end{enumerate}

	\noindent \textbf{Выводы:}
	
	В ходе лабораторной работы изучена система команд и способы адресации процессоров архитектуры x86. Проанализировано размещение глобальных данных в памяти, составлена карта памяти. Установлено, что при обращении к элементам массива по индексу используется многокомпонентная масштабно-индексная адресация с масштабом 4 (размер типа \texttt{int}). Замена индексации на указатель упрощает команду записи до короткой косвенно-регистровой (\texttt{mov [rcx], eax}) без вычисления адреса на каждой итерации. Проверка условий беззнаковых переменных (\texttt{unsigned}) транслируется в беззнаковые переходы (\texttt{jb}, \texttt{jae}), а умножение на степень двойки ($i \cdot 8$) компилятор эффективно заменяет побитовым сдвигом влево (\texttt{shl eax, 3}).

	%КАРТИНКИ
	%\begin{figure}[H]
	%	\centering
	%	\includegraphics[width=0.85\textwidth]{install_deb.png}
	%	\caption{Начало установки гостевой ОС Debian 12}
	%	\label{fig:vbox_setup}
	%\end{figure}

	%---------------------------------------------------------------------------------
	%ССЫЛКИ НА САЙТ

	%\href{https://www.virtualbox.org/wiki/Downloads}{официального сайта VirtualBox}

	%---------------------------------------------------------------------------------
	%ССЫЛКА НА ПРИЛОЖЕНИЕ

	%\texttt{zlib}\cite{zlib_docs}

	%---------------------------------------------------------------------------------
	%ССЫЛКА НА РИСУНОК

	%показано на рис. \ref{fig:vbox_setup}

	%---------------------------------------------------------------------------------
	%ТАБЛИЦА С ПЕРЕВОДОМ НА СЛЕДУЮЩУЮ СТРАНИЦУ

	%\begin{longtable}{|l|p{0.75\textwidth}|}
	%	\caption{Функции библиотеки ncurses, используемые в TUI} \label{tab:ncurses_funcs} \\
	%	\hline
	%	\textbf{Функция} & \textbf{Назначение и описание алгоритма работы} \\ \hline
	%	\endhead
	%	\texttt{initscr()} & Инициализирует библиотеку, выделяет память под внутренние структуры окна и переводит терминал ОС в графический режим. \\ \hline
	%	\texttt{cbreak()} & Отключает строковую буферизацию терминала, позволяя программе мгновенно реагировать на нажатия клавиш без ожидания \texttt{Enter}. \\ \hline
	%	\texttt{noecho()} & Отключает автоматический вывод (эхо) вводимых пользователем символов на экран, позволяя программе самой управлять отрисовкой. \\ \hline
	%	\texttt{clear()} & Очищает виртуальный экран для подготовки к отрисовке нового кадра меню. \\ \hline
	%	\texttt{move(y, x)} & Осуществляет абсолютное позиционирование курсора по координатам \texttt{Y} (строка) и \texttt{X} (столбец). \\ \hline
	%	\texttt{printw()} & Аналог функции \texttt{printf} для \texttt{ncurses}: выводит отформатированный текст в текущей позиции курсора. \\ \hline
	%	\texttt{getch()} & Ожидает и считывает код нажатой пользователем клавиши (включая стрелки навигации и функциональные клавиши). \\ \hline
	%	\texttt{getnstr()} & Безопасно считывает строку от пользователя с аппаратным ограничением максимального количества символов (защита от переполнения). \\ \hline
	%	\texttt{scanw()} & Аналог функции \texttt{scanf}, считывает форматированный ввод пользователя (применяется для ввода числовых данных). \\ \hline
	%	\texttt{endwin()} & Выгружает библиотеку из памяти, восстанавливает историю команд и возвращает терминал в стандартный режим работы ОС. \\ \hline
	%\end{longtable}
	%---------------------------------------------------------------------------------

\end{document}
